% Source: https://artofproblemsolving.com/community/c3h1435307_statistics_basic_vocabulary
% Stable identity: aops-topic-1435307-q01
\textbf{(a)} The sum is $3+5+6+6+8+10+11+15=64$, so the mean is $\dfrac{64}{8}=8$. With $8$ values (already in order) the median is the average of the 4th and 5th values: $\dfrac{6+8}{2}=7$. The mode (most frequent value) is $6$. The range is $15-3=12$.

\textbf{(b)} The lower half is $3,5,6,6$, with median $Q_1=\dfrac{5+6}{2}=5.5$; the upper half is $8,10,11,15$, with median $Q_3=\dfrac{10+11}{2}=10.5$. The five-number summary is
\[
3,\quad 5.5,\quad 7,\quad 10.5,\quad 15 .
\]

\textbf{(c)} $\text{IQR}=10.5-5.5=5$; the middle half of the data lies in an interval of length $5$.

\textbf{(d)} The distances from the mean $8$ are $5,3,2,2,0,2,3,7$, with sum $24$, so the MAD is $\dfrac{24}{8}=3$.

The box plot has a box from $5.5$ to $10.5$ with a line inside it at the median $7$, and whiskers from the box out to the minimum $3$ and the maximum $15$.

So: mean $8$, median $7$, mode $6$, range $12$; five-number summary $3,\ 5.5,\ 7,\ 10.5,\ 15$; $\text{IQR}=5$; $\text{MAD}=3$.
